\chapter{Elliptische Kurven}
Dieses Kapitel legt die Grundlagen für die Arbeit mit elliptischen Kurven. Neben der allgemeinen Gleichung wird auch die Gruppenstruktur eingeführt. Ein großer Punkt spielt dabei der Beweis der Assoziativität. Außerdem wird auch die projektive Ebene eingeführt

\section{Grundlagen}
In diesem Abschnitt werden die Grundlagen für elliptische Kurven eingeführt, auf denen anschließend die Beschreibung
der Gruppenstruktur aufbaut. Die Darstellung in diesem Abschnitt orientiert sich an \cite[S.~9~ff.]{washington08}, \cite[S.~1, 6]{plaumann20} und für die Diskriminante an \cite[S.~470~f.]{karpf24}.

Beginnen wir also mit der Definition einer Elliptische Kurve.
\begin{definition}
    Eine elliptische Kurve über einem Körper K wird definiert mit
    \[
        y^2+a_1xy+a_3y=x^3+a_2x^2+a_4x+a_6,
    \]
    wobei wir das die Weierstraß-Gleichung nennen.
    Wenn nun gilt $\operatorname{char}(K)\neq 2,3$, so können wir die kurze
    Weierstraß-Gleichung nutzen
    \[
        y^2=x^3+Ax+B.
    \]
    Wobei A und B konstante Werte sind.
    Dabei sind $A,B\in K$ und es gilt
    \[
        4A^3+27B^2\neq 0.
    \]
\end{definition}
Der Punkt im Unendlichen wird mit \(\infty =\mathcal{O}\) bezeichnet. Für einen Erweiterungskörper $L$ von $K$, also $K\subseteq L$, bezeichnet
\[
    E(L)=\{\mathcal{O} \} \cup \{(x,y)\in L\times L\mid y^2=x^3+Ax+B\}
\]
die Menge aller rationalen Punkte der elliptischen Kurve $E$ und den Punkt im Unendlichen. Nach Definition ist der Punkt $\mathcal{O}$ immer enthalten, warum dieser benötigt wird, wird im nächsten Abschnitt ausgeführt. Nicht über jedem Körper lassen sich Elliptische Kurven sinnvoll optisch darstellen, daher wird zuerst der Körper der reellen Zahlen betrachtet. Hierbei gibt es zwei grundlegende Formen.

Über dem Körper \(\mathbb{R}\) kann die Kurve je nach Anzahl der reellen Nullstellen verschiedene Formen annehmemen. Zum Beispiel hat $y^2 = x^3 - 3x$ drei verschiedene reelle Nullstellen, nämlich $x = 0$ und $x = \pm \sqrt{3}$. Im Gegensatz dazu hat die Kurve $y^2 = x^3 + 3x$ lediglich eine reelle Nullstelle bei $x = 0$.
\begin{figure}[H]
    \centering
    \includegraphics[width=0.35\linewidth]{kap2_bs1.png}
    \hfill
    \includegraphics[width=0.35\linewidth]{kap2_bs2.png}
    \caption{Elliptische Kurven über $\mathbb{R}$:
        links $y^2=x^3-3x$, rechts $y^2=x^3+3x$.
        Eigene Darstellung mit GeoGebra.}
    \label{fig:vergleich}
\end{figure}
Wie man in der Darstellung sehen kann, ändert sich die Form der Kurve über dem Körper $\mathbb{R}$, wenn mehrere reelle Nullstellen exestieren. Das Polynom \(x^3+Ax+B\) darf daher keine mehrfachen Nullstellen besitzen, dies ist äquivalent dazu, dass die Kurve \emph{nicht singulär} ist.


\begin{definition}
    Sei $f \in K[x,y]$ und
    \[
        E = V(f) = \{(x,y)\in K^2 \mid f(x,y)=0\}
    \]
    eine affine ebene Kurve. Ein Punkt \((a,b)\in E\) heißt \textbf{singulär}, wenn,
    \[
        f_x(a,b)=0 \quad \text{und}\quad f_y(a,b)=0.
    \]
    Andernfalls heißt der Punkt regulär.

    Die Kurve \(E\) heißt \textbf{singulär}, wenn sie einen singulären Punkt besitzt.
\end{definition}

Die elliptische Kurve soll diese Eigenschaft erfüllen, daher wird untersucht welche Bedingung dafür gelten muss. Dafür wird die kruze Weierstraßgleichung einer Kurve $E$ umgeschrieben zu
\[
    f(x,y)=y^2-x^3-Ax-B
\]
dann sind die Ableitung
\[
    f_x= -3x^2-A \quad \text{ und } \quad f_y=2y
\]
Es muss nun für eine der beiden pratiellen Ableitungen ungleich 0 gelten,
\[
    f_x= -3x^2-A = 0 \Leftrightarrow 3x^2+A = 0
\]
und
\[
    f_y=2y = 0 \Leftrightarrow y = 0, \quad \operatorname{char}(K)\neq 2
\]
daraus folgt ein singulärer Punkt muss die Form $(x,0)$ haben und zusätzlich muss \(x^3+Ax+B = 0\) gelten, sodass der Punkt auf der Kurve liegt. Also x darf nicht eine Nullstelle von \(x^3+Ax+B\) und der Ableitung \(3x^2+A\) sein. Die Diskriminante für ein kubisches Polynom
\[
    p(x)=x^3+Ax+B
\]
ist hierbei
\[
    \Delta  = -4A^3-27B^2.
\]
Damit besitzt $p$ genau dann keine mehfachen Nullstellen, wenn
\[
    4A^3+27B^2 \neq 0.
\]
Diese Bedingung ist äquivalent mit der Zusatzbedingung für $A,B\in K$ bei der Definition einer elliptischen Kurve.

Für die Nullstellen $r_1, r_2, r_3$ des Polynoms gilt
\[
    ((r_1 - r_2)(r_1 - r_3)(r_2 - r_3))^2 = -(4A^3 + 27B^2).
\]
Insbesondere folgt daraus, dass die Nullstellen genau dann paarweise verschieden sind, wenn die Diskriminante $\Delta = -(4A^3 + 27B^2)$ ungleich Null ist. Im vorliegenden Fall ist $\Delta=108 > 0$, sodass drei voneinader verschiedene reelle Nullstellen vorliegen. Das Polynom $x^3+3x$ besitzt dagegen wegen $\Delta=-108<0$ genau eine reelle Nullstelle. In beiden Fällen ist die Diskriminante ungleich null, sodass die zugehörigen Weierstraßkurven nicht singulär sind.

\section{Gruppenstruktur}
In diesem Abschnitt wird die Gruppenoperation,also die Addition zweier Punkte auf einer elliptischen Kurve eingeführt. Zunächst wird die geometrische Konstruktion der Punktaddition über $\mathbb{R}$ erläutert. Anschließend werden daraus die Additionsformeln hergeleitet und die Gruppeneigenschaften untersucht. Die Darstellung orientiert sich an \cite[S.~12--15]{washington08}, \cite[S.~8--19,30--31]{silverman06}, \cite[S.~179~ff.]{kapl05}
und \cite[S.~99]{buell24}. Der Beweis der Wohldefiniertheit folgt \cite[S.~477--481]{zwegers24}.

\subsection{Geometrische Konstruktion der Punktaddition}
Im Folgenden sei
\[
    E:\quad y^2=x^3+Ax+B
\]
eine elliptische Kurve über einem Körper $K$ mit $\operatorname{char}(K)\neq 2,3$. Zunächst wird die Addition zweier Punkte betrachtet. Seien dazu
\[
    P=(x_1,y_1),\qquad Q=(x_2,y_2)\in E(K).
\]
Gesucht ist der Punkt \(R=P+Q\). Genauer handelt es sich bei $+$ um die auf der elliptischen Kurve definierte Gruppenoperation. Zur Vereinfachung der Notation wird im Folgenden jedoch auf die ausführlichere Schreibweise $P+_E Q$ verzichtet. Zur Veranschaulichung der geometrischen Konstruktion wird zunächst ein Beispiel über den reellen Zahlen betrachtet. Hierzu sei
\[
    E:\quad y^2=x^3-4x+10.
\]
Im ersten Schritt wird die Gerade $L$ durch $P$ und $Q$ gelegt.
\begin{figure}[htbp]
    \centering
    \includegraphics[width=0.5\linewidth]{kap2_2_bsp1.png}
    \caption{Gerade $L$ durch die Punkte $P$ und $Q$}
    \label{fig:punktaddition-gerade}
\end{figure}
Im zweiten Schritt wird der Schnittpunkt $R'$ von der Gerade $L$ und $E$ mit der elliptischen Kurve $E$ betrachtet. Dieser wird an der x-Achse gespiegelt. Der dadurch entstandene Punkt wird mit $R$ bezeichnet und als $R=P+Q$ definiert. Dabei ist zubeachten, dass es sich bei der Addition zweier Punkte auf einer elliptischen Kurven nicht um die koordinatenweise Punktaddition handelt.

\begin{figure}[htbp]
    \centering
    \includegraphics[width=0.5\linewidth]{kap2_2_bsp2.png}
    \caption{Dritter Schnittpunkt $R'$ der Geraden $L$ mit der Kurve
        $E$ und Spiegelung an der $x$-Achse}
    \label{fig:punktaddition-dritter-sp}
\end{figure}

Als nächstes wird der Fall betrachtet, wenn ein Punkt $P$ auf der Kurve $E$ auf sich selbst addiert wird. Da in diesem Fall keine Gerade durch zwei verschiedene Punkte $P$ und $Q$ bestimmt werden kann, wird der Fall betrachtet, dass sich $Q$ und $P$ annähern. Die Sekante durch $P$ und $Q$ geht dabei in die Tangente an $E$ im Punkt $P$ über.
\begin{figure}[htbp]
    \centering
    \includegraphics[width=0.5\linewidth]{kap2_2_bsp3.png}
    \caption{Tangente an die Kurve $E$ im Punkt $P$}
    \label{fig:punktaddition-spiegelung}
\end{figure}

Der weitere Ablauf entspricht der zuvor beschriebenen Punktaddition. Die Tangente schneidet die Kurve $E$ in einem weiteren Punkt $R'$.
Durch Spiegelung dieses Punktes an der $x$-Achse entsteht der Punkt $R=P+P=2P$.
\begin{figure}[htbp]
    \centering
    \includegraphics[width=0.5\linewidth]{kap2_2_bsp4.png}
    \caption{Spiegelung von $R'$ und Konstruktion des Punktes $R=2P$}
    \label{fig:punktverdoppelung-tangente}
\end{figure}
Abschließend wird der Fall betrachtet, wenn $L$ keinen Schnittpunkt mit $E$ hat. Es wird also die Addition der Punkte $P$ und $P'=-P$ betrachtet. Der Punkt $-P$ entsteht durch Spiegelung von $P$ an der $x$-Achse. Die Gerade
durch $P$ und $-P$ ist daher vertikal.
\begin{figure}[htbp]
    \centering
    \includegraphics[width=0.5\linewidth]{kap2_2_bsp5.png}
    \caption{Schnittpunk von L im Punkt $\mathcal{O}$}
    \label{fig:punktverdopplung-ergebnis}
\end{figure}
Für diesen Fall wird der zusätzliche Punkt $\mathcal{O}$ genutzt, der als Punkt im Unendlichen bezeichnet wird. Jede vertikale Gerade wird so aufgefasst, dass sie die elliptische Kurve zusätzlich im Punkt $\mathcal O$ schneidet. Damit wird
\[
    P+(-P)=\mathcal O
\]
definiert. Der Punkt $\mathcal O$ bildet später das neutrale Element der Gruppenstruktur. Eine genauere Beschreibung mithilfe projektiver Koordinaten erfolgt in Abschnitt \ref{subsec:punkt-im-unendlichen}.

\subsection{Herleitung der Additionsformeln}
Seien $P_1=(x_1,y_1)$ und $P_2=(x_2,y_2)$ Punkte auf $E$ mit $P_1,P_2\neq \mathcal{O}$. Zuerst wird der Fall $P\neq Q$ und \(x_1\neq x_2\) betrachtet. Dann hat die Gerade \(L\) durch durch \(P_1\) und \(P_2\) die Form
\[
    y=m(x-x_1)+y_1
    \qquad \text{mit} \qquad
    m=\frac{y_2-y_1}{x_2-x_1}.
\]
Zur Bestimmung des dritten Schnittpunkts wird die Geradengleichung in die Kurvengleichung eingesetzt
\[
    E: y^2=x^3+Ax+B
\]
ein. Es gilt also
\[
    \bigl(m(x-x_1)+y_1\bigr)^2=x^3+Ax+B.
\]
Nach Ausmultiplizieren und Umstellen erhält man
\[
    0=x^3-m^2x^2+\bigl(2m^2x_1-2my_1+A\bigr)x
    +\bigl(B-m^2x_1^2+2mx_1y_1-y_1^2\bigr).
\]
Die Nullstellen des kubischen Polynoms sind die \(x\)-Koordinaten der drei Schnittpunkte von \(L\) mit \(E\), also \(x_1\), \(x_2\) und \(x_3'\). Daher gilt
\[
    0=(x-x_1)(x-x_2)(x-x_3')
    =x^3-(x_1+x_2+x_3')x^2+(x_1x_2+x_1x_3'+x_2x_3')x-x_1x_2x_3'
\]
Mit Kooeffizientenvergleich erhält man
\[
    x_1+x_2+x_3' = m^2 \Leftrightarrow  x_3' = m^2-x_1-x_2
\]
Da \(P_3'=(x_3',y_3')\) auf der Geraden \(L\) liegt, gilt
\[
    y_3'=m(x_3'-x_1)+y_1.
\]
Die Summe \(P_1+P_2\) ist die Spiegelung von \(P_3'\) an der \(x\)-Achse. Somit ist
\[
    P_1+P_2=P_3=(x_3,y_3)
\]
mit
\[
    x_3=x_3'=m^2-x_1-x_2,
    \qquad
    y_3=-y_3'=m(x_1-x_3)-y_1.
\]
Ist $x_1=x_2$ und $y_1\neq y_2$, so gilt wegen $y_1^2=y_2^2$ notwendigerweise $y_2=-y_1$. Die Gerade durch $P$ und $Q$ ist somit vertikal und schneidet die projektive Kurve im Punkt $\mathcal{O}$. Daher wird
\[
    P_1+P_2=\mathcal{O}
\]
definiert. Wenn man nun den Punkt an der x-Achse spiegelt erhält man den gleichen Punkt in $\mathcal{O}$, daher gilt für $P_1+P_2=\mathcal{O}$. Jetzt betrachtet man den Fall $P_1=P_2=(x_1,y_1)$. Wenn zwei Punkte auf einer Kurve sehr nah beinander sind, dann approximiert die Gerade eine Tangente. Daher wird die Gerade $L$ durch die beiden Punkte als Tangente betrachten. Mit Hilfe der Ableitung kann man die Steigungn $m$ von der Gerade $L$ bestimmen
\[
    2y\frac{dy}{dx}=3x^2+A \Leftrightarrow \quad m=\frac{dy}{dx}=\frac{3x_1^2+A}{2y_1}.
\]
Für den Fall $y_1=0$ ist die Gerade vertikal und es gilt wieder $P_1+P_2=\mathcal{O}$. Daher wird der Fall $y_1\neq 0$ betrachtet. Die Gleichung der Gerade ist in diesem Fall
\[
    y=m(x-x_1)+y_1.
\]
Setzt man nun die Gerade in die Kurvengleichung ein, erhält man
\begin{align}
     & (m(x-x_1)+y_1)^2=x^3+Ax+B                                 \\
     & 0=x^3-m^2x^2+(2m^2x_1-2my_1+A)x+B-m^2x_1^2+2my_1x_1-y_1^2
\end{align}
Man kennt von dieser Gleichung eine Nullstelle $x_1$, da es sich bei $L$ um eine Tangente von $E$ handelt ist $x_1$ eine doppelte Nullstelle. Daher kann man analog zu vorher
\[
    x_3=m^2-2x_1,\quad y_3=m(x_1-x_3)-y_1
\]
rechnen. Für den Punkt im Unendlichen wird die Addition durch
\[
    P+\mathcal{O}=\mathcal{O}+P=P
\]
definiert. Dies entspricht geometrisch der Tatsache, dass die vertikale
Gerade durch $P$ den zweiten affinen Schnittpunkt $-P$ besitzt und
zusätzlich durch $\mathcal{O}$ verläuft. Nach Spiegelung von $-P$
an der $x$-Achse wird wieder $P$ erhalten. Diese Aussage kann erweitert werden, sodass \(\mathcal{O}+\mathcal{O}=\mathcal{O}\) inkludiert ist.
\begin{definition}[Punktaddition]
    \label{def:punktaddition}
    Die Punktaddition für einen Punkt $P_1+P_2=P_3=(x_3,y_3)$ definieren wir wie folgt:
    \begin{enumerate}
        \item Wenn $x_1\neq x_2$, dann
              \[x_3=m^2-x_1-x_2, \quad y_3=m(x_1-x_3)-y_1 \quad \text{ mit  } m=\frac{y_2-y_1}{x_2-x_1}.\]
        \item Wenn $x_1=x_2$, aber $y_1\neq y_2$, dann ist
              \[P_1+P_2=\mathcal{O}.\]
        \item Wenn $P_1=P_2$ und $y_1\neq 0$, dann
              \[x_3=m^2-2x_1, \quad y_3=m(x_1-x_3)-y_1, \quad \text{ mit } m=\frac{3x_1^2+A}{2y_1}.\]
        \item Wenn $P_1=P_2$ und $y_1=0$ dann,
              \[P_1+P_2=\mathcal{O}.\]
    \end{enumerate}

    Zusätzlich wird festgelegt, dass
    \[
        P+\mathcal{O}=\mathcal{O}+P=P
    \]
    für alle Punkte $P$ auf $E$.
\end{definition}

\begin{lemma}[Wohldefiniertheit der Punktadditon]
    \label{lem:wohldefiniertheit}
    Seien \(P_1=(x_1,y_1)\) und \(P_2=(x_2,y_2)\) Punkte aus $E\setminus \{\mathcal{O}\}$ mit \((x_2,y_2)\neq (x_1,-y_1)\). Dann gilt
    \[
        x_3 = m_{12}^2-x_1-x_2 \quad \text{ und } \quad y_3=-y_1-m_{12}(x_3-x_1)
    \]
    mit
    \[
        m_{12}=m(P_1,P_2)=
        \begin{cases}
            \dfrac{y_2-y_1}{x_2-x_1}, & \text{wenn } x_1 \neq x_2, \\[6pt]
            \dfrac{3x_1^2+A}{2y_1},   & \text{wenn } P_1=P_2.
        \end{cases}
    \]
    Dann ist \(m_{12}\) wohldefiniert und es gilt
    \[
        (x_3,y_3)\in E(K).
    \]
\end{lemma}

\begin{proof}
    Zunächst ist zu zeigen, dass die angegebene Definition von \(m_{12}\) in den angegebenen Fällen wohldefiniert ist. Für \(x_1\neq x_2\) ist der Nenner \(x_2-x_1\) von null verschieden, damit ist die Steigung der Sekante durch \(P_1\) und \(P_2\) wohldefiniert und es gilt
    \[
        m_{12}=\frac{y_2-y_1}{x_2-x_1}\in K.
    \]
    Sei nun \(x_1=x_2\). Da \(P_1,P_2\in E(K)\) liegen, gilt mit \(f(x)=x^3+Ax+B\)
    zunächst
    \[
        y_1^2=f(x_1)=f(x_2)=y_2^2.
    \]
    Daraus folgt \(y_2=y_1\text{ oder } y_2=-y_1\). Es kann aufgrund der Vorraussetzung des Lemmas nicht \(y_2=-y_1\) gelten, daher muss $y_2=y_1$ gelten und damit auch \(P_1=P_2\). Angenommen, es gilt\(y_1=0\), dann gilt $y_2=y_1=0=-y_1$ und somit im Widerspruch zur Vorraussetzung. Daher muss \(y_1\neq 0\) gelten. Wegen \(\operatorname{char}(K)\neq2\) ist daher \(2y_1\neq0\), sodass auch
    \[
        m_{12}=\frac{3x_1^2+A}{2y_1}\in K
    \]
    wohldefiniert ist. Zur Bestimmung des summierten Punktes wird das Polynom
    \[
        g(x)=f(x)-(y_1+m_{12}(x-x_1))^2\in K[x].
    \]
    eingeführt. Die Nullstellen von \(g\) entsprechen den \(x\)-Koordinaten der Schnittpunkte der Geraden \(y=y_1+m_{12}(x-x_1)\) mit der elliptischen Kurve. Da \(P_1\in E(K)\) gilt,
    \[
        g(x_1)=f(x_1)-\left(y_1+m_{12}(x_1-x_1)\right)^2=f(x_1)-y_1^2=0.
    \]
    Somit ist \(x_1\) eine Nullstelle von \(g\). Für \(x_1\neq x_2\) gilt umgestellt \( y_1+m_{12}(x_2-x_1)=y_2\). Daher folgt
    \[
        g(x_2)=f(x_2)-\left(y_1+m_{12}(x_2-x_1)\right)^2 =f(x_2)-y_2^2=0,
    \]
    sodass auch  $x_2$ eine Nullstelle von \(g\) ist.

    Für den Fall \(x_1=x_2\), entspricht die Gerade der Tangente an \(E\) im Punkt \(P_1\). Es ist somit zu zeigen, dass \(x_1\) eine doppelte Nullstelle von \(g\) ist. Die Ableitung von \(g\) lautet
    \[
        g'(x)=f'(x)-2m_{12}\left(y_1+m_{12}(x-x_1)\right) \in K[x].
    \]
    Für $x_1$ gilt dann
    \[
        g'(x_1)=f'(x_1)-2m_{12}(y_1)= (3x_1^2+A)-2\frac{3x_1^2+A}{2y_1}y_1=0.
    \]
    Damit ist \(x_1\) eine doppelte Nullste von \(g\). Es lässt sich somit \(g\) in beiden Fällen in der Form
    \[
        g(x)=(x-x_1)(x-x_2)(x-x_3)
    \]
    schreiben, wobei im Fall \(x_1=x_2\) der Fakor \((x-x_1)\) doppelt auftritt. Da \(g\) ein normiertes Polynom dritten Grades ist, kann \(x_3\in K\) durch einen Koeffizientenvergleich bestimmt werden. Es gilt
    \begin{equation}
        \label{eq:faktorisierung-wohldefiniertheit}
        f(x)-\left(y_1+m_{12}(x-x_1)\right)^2
        =(x-x_1)(x-x_2)(x-x_3).
    \end{equation}
    Der Koeffizient von \(x^2\) auf der linken Seite ist \(-m_{12}^2\) während der Koeffizient auf der rechten Seite \(-(x_1+x_2+x_3)\) beträgt. Daraus folgt \(-m_{12}^2=-(x_1+x_2+x_0) \)
    und somit \(x_3=m_{12}^2-x_1-x_2\). Da $x_3$ eine Nullstelle von $g$ ist, gilt
    \[
        g(x_3)
        =
        f(x_3)
        -
        \left(y_1+m_{12}(x_3-x_1)\right)^2
        =0.
    \]
    Folglich ist
    \[
        f(x_3)
        =
        \left(y_1+m_{12}(x_3-x_1)\right)^2.
    \]

    Nach Definition der Additionsformel ist
    \[
        y_3
        =
        -y_1-m_{12}(x_3-x_1)
        =
        -\left(y_1+m_{12}(x_3-x_1)\right).
    \]
    Durch Quadrieren folgt
    \[
        y_3^2
        =
        \left(y_1+m_{12}(x_3-x_1)\right)^2.
    \]
    Damit erfüllt \((x_3,y_3)\) die Kurvengleichung
    \[
        y_3^2=x_3^3+Ax_3+B
    \]
    und somit gilt \((x_3,y_3)\in E(K)\).
\end{proof}
Damit ist die oben definierte Punktaddition abgeschlossen auf \(E\).
\begin{bemerkung}
    \label{bem:sum-ist-identitaet}
    Damit ist der nichttriviale Fall der Punktaddition wohldefiniert.
    Für $P_2=-P_1$ gilt nach Definition
    \[
        P_1+P_2=\mathcal{O}\in E,
    \]
    und für $P\in E$ gilt
    \[
        P+\mathcal{O}=\mathcal{O}+P=P\in E.
    \]
    Somit ist die Punktaddition insgesamt abgeschlossen auf $E$.
\end{bemerkung}

\begin{theorem}
    \label{theo:eigenschaften-addition}
    Die Menge $E$ bildet bezüglich der Punktaddition eine abelsche Gruppe
    mit neutralem Element $\mathcal{O}$. Insbesondere gelten für alle
    $P_1,P_2,P_3\in E$:
    \begin{enumerate}
        \item[(i)] \textbf{Neutrales Element:}
              \(P+\mathcal{O}=\mathcal{O}+P=P\).

        \item[(ii)] \textbf{Inverses Element:}
              Zu jedem $P\in E$ existiert $-P\in E$ mit
              \(P+(-P)=\mathcal{O}\).

        \item[(iii)] \textbf{Assoziativität:}
              \((P_1+P_2)+P_3=P_1+(P_2+P_3)\).

        \item[(iv)] \textbf{Kommutativität:}
              \(P_1+P_2=P_2+P_1\).
    \end{enumerate}
\end{theorem}

\begin{bemerkung}
    Sei $E$ eine elliptische Kurve über $\mathbb{R}$ und $P=(x,y)\in E$, dann ist das Inverse Element von $P$ definiert als $-P=(x,-y)$ und somit symmetrisch zu $P$ an der $x$-Achse.
\end{bemerkung}

\begin{proof}
    Die Existenz des neutralen Elements folgt unmittelbar aus
    Definition~\ref{def:punktaddition}, da für alle $P\in E$
    \[
        P+\mathcal{O}=\mathcal{O}+P=P
    \]
    gilt.

    Die Kommutativität folgt aus der Definition der Punktaddition. Es werden die in Definition~\ref{def:punktaddition} auftretenden Fälle unterschieden.
    \begin{enumerate}
        \item Sei $x_1\neq x_2$. Dann gilt
              \[
                  m_{12}
                  =\frac{y_2-y_1}{x_2-x_1}
                  =\frac{y_1-y_2}{x_1-x_2}
                  =m_{21}.
              \]
              Damit stimmen die $x$-Koordinaten der Summen überein
              \[
                  m_{12}^2-x_1-x_2
                  =
                  m_{21}^2-x_2-x_1.
              \]
              Für die $y$-Koordinaten gilt
              \[
                  \begin{aligned}
                      y_3'-y_3
                       & =
                      -y_2-m_{21}(x_3-x_2)
                      +y_1+m_{12}(x_3-x_1)    \\
                       & =
                      y_1-y_2+m_{12}(x_2-x_1) \\
                       & =0.
                  \end{aligned}
              \]
              Somit folgt $P_1+P_2=P_2+P_1$.

        \item Sei $x_1=x_2$ und $y_1\neq y_2$. Dann gilt nach Definition
              \[
                  P_1+P_2=\mathcal{O}=P_2+P_1.
              \]

        \item Sei $P_1=P_2$ und $y_1\neq 0$. Da beide Summanden identisch sind,
              ergibt die Tangentenformel unabhängig von der Reihenfolge dieselbe
              Steigung und damit denselben Summenpunkt. Also gilt
              \[
                  P_1+P_2=P_2+P_1.
              \]

        \item Sei $P_1=P_2$ und $y_1=0$. Dann ist $P_1=-P_1$ und nach Definition gilt
              \[
                  P_1+P_2=\mathcal{O}=P_2+P_1.
              \]
    \end{enumerate}
    Für $P=(x_1,y_1)\in E$ sei
    \[
        -P=(x_1,-y_1).
    \]
    Da
    \[
        (-y_1)^2=y_1^2,
    \]
    liegt auch $-P$ auf $E$. Nach Definition der Punktaddition gilt
    \[
        P+(-P)=\mathcal{O}.
    \]
    Für $P=\mathcal{O}$ gilt $-\mathcal{O}=\mathcal{O}$.
    Damit besitzt jedes Element ein inverses Element.

    Die Assoziativität wird in Abschnitt~\ref{sec:assoziativitaet}
    gesondert bewiesen.
\end{proof}

\begin{beispiel}
    Sei $E$ eine elliptische Kurve mit
    \[
        y^2=x^3-4x+4
    \]
    als Gleichung. Auf dieser Kurve liegen die Punkte $P_1=(0,2)$ und $P_2=(2,2)$. Nun soll die Summe aus $P_1$ und $P_2$ mit Hilfe der Additionsformel berechnet werden. Es ist $x_1\neq x_2$ und daher folgt
    \[
        m = \frac{2-2}{2-0}=0.
    \]
    Damit ergibt sich
    \[
        x_3 = 0^2-0-2=-2 \quad \text{und} \quad y_3 = 0\cdot(0-(-2))-2 =-2.
    \]
    Somit folgt
    \[
        P_1+P_2=(0,2)+(2,2)= (-2,-2)
    \]
    Nach demselben Vorgehen können weitere Punkte auf $E$ miteinander addiert werden. Beispielsweise gilt
    \[
        (0,2)+(-2,-2)=(6,-14).
    \]
\end{beispiel}

\begin{beispiel}
    Als Folge ist es der algebraischen Additionsformeln können auch Punkte auf elliptischen Kurven über Körpern addiert werden, für die keine vergleichbare geometrische Darstellung wie über $\mathbb{R}$ existiert. Sei dazu $E$ die elliptische Kurve über dem endlichen Körper $\mathbb{F}_7$ mit \(E: \quad y^2=x^3-4x+4\). Zunächst ist zu überprüfen, dass die Kurve nicht singulär ist. Mit $A=-4$ und $B=4$ gilt
    \[
        4A^3+27B^2
        =
        4\cdot(-4)^3+27\cdot4^2
        =
        176
        \equiv1\bmod 7.
    \]
    Da
    \[
        4A^3+27B^2\not\equiv0\bmod7,
    \]
    ist $E$ eine elliptische Kurve über $\mathbb{F}_7$.
    Nun seien \(P_1=(0,2)\text{ und }P_2=(5,5)\) zwei Punkte aus $E(\mathbb{F}_7)$. Da $x_1\neq x_2$ gilt, folgt mit der Additionsformel
    \[
        m=\frac{5-2}{5-0}=\frac{3}{5} = 3\cdot 5^{-1}\equiv 3\cdot 3 7= 9 = 2 \bmod 7
    \]
    und damit
    \[
        x_3=2^2-0-5=-1\equiv 6 \bmod 7\quad \text{und}\quad  y_3=2(0-6)-2=-14 \equiv 0\bmod 7.
    \]
    Somit folgt
    \[
        P_1+P_2=(6,0).
    \]
\end{beispiel}

\subsection{Skalar mal Punkt}
Die Punktaddition ermöglicht die Definition ganzzahlig Vielfacher eines Kurvenpunkts. Für $P\in E(K)$ und
$m\in\mathbb{N}_{>0}$ ist
\[
    mP=\underbrace{P+\cdots+P}_{m\text{-mal}}.
\]
Zusätzlich gelten
\[
    0P=\mathcal{O}
    \qquad\text{und}\qquad
    (-m)P=-(mP).
\]
Diese Operation wird als Skalarmultiplikation bezeichnet.

Mithilfe der Binärdarstellung von $m$ lässt sich $mP$ durch wiederholtes Verdoppeln und Addieren mit $\mathrm{O}(\log m)$ Gruppenoperationen berechnen. Das zugehörige Double-and-Add-Verfahren wird in Abschnitt~\ref{sec:double-and-add} erläutert.

\section{Projektive Ebene und der Punkt bei Unendlich}
\label{subsec:punkt-im-unendlichen}
Der projektive Raum dient dazu, den Punkt im Unendlichen $\mathcal{O}$ genauer zu definieren, da sich dieser in affinen Koordinaten nicht darstellen lässt.  Um diesen Punkt geometrisch zu beschreiben, wird die
affine Ebene in die projektive Ebene eingebettet. Das Kapitel basiert auf \cite[S.~18--20]{washington08}, \cite[S.~245--246]{clader25}, \cite[S.~200--204]{forster15}, \cite[S.~15--22]{werner02} und \cite[S.~9]{liebscher17}.

\begin{definition}[Affine Ebene]
    Sei $K$ ein Körper. Dann ist die affine Ebene über $K$ definiert durch
    \[
        \mathbb{A}^2_K
        =
        K \times K
        =
        \{(x,y) \mid x,y \in K\}.
    \]
\end{definition}
In affinen Koordinaten kann eine elliptische Kurve beispielsweise durch eine Gleichung der Form
\[
    y^2 = x^3 + Ax + B
\]
beschrieben werden. Der Punkt im Unendlichen $\mathcal{O}$ lässt sich jedoch nicht als Punkt $(x,y)\in \mathbb{A}^2_K$ darstellen. Aus diesem Grund wird die projektive Ebene eingeführt.

\begin{definition}[Projektive Ebene]
    Sei $K$ ein Körper. Die projektive Ebene über $K$ ist definiert als
    \[
        \mathbb{P}^2_K
        =
        \left(K^3 \setminus \{(0,0,0)\}\right)/\sim,
    \]
    wobei für
    \[
        (X,Y,Z),(X',Y',Z') \in K^3 \setminus \{(0,0,0)\}
    \]
    die Äquivalenzrelation
    \[
        (X,Y,Z)\sim(X',Y',Z')
    \]
    genau dann gilt, wenn ein $\lambda \in K^\times$ existiert mit
    \[
        (X,Y,Z)
        =
        \lambda(X',Y',Z').
    \]
    Die Äquivalenzklasse eines Tripels $(X,Y,Z)$ wird mit
    \[
        [X:Y:Z]
    \]
    bezeichnet und heißt Punkt der projektiven Ebene.
\end{definition}

Die Koordinaten $[X:Y:Z]$ heißen homogene Koordinaten. Für jedes $\lambda\in K^\times$ gilt daher
\[
    [X:Y:Z]
    =
    [\lambda X:\lambda Y:\lambda Z].
\]
Jedes Tripel
\[
    (X,Y,Z)\in K^3\setminus\{(0,0,0)\}
\]
bestimmt somit einen Punkt in $\mathbb{P}^2_K$. Das Tripel $(0,0,0)$ ist ausgeschlossen und definiert keinen projektiven Punkt. Die affine Ebene lässt sich mit dem Teil der projektiven Ebene identifizieren, für den $Z\neq 0$ gilt.  Wegen
\[
    [X:Y:Z]=\left[\frac{X}{Z}:\frac{Y}{Z}:1\right]
\]
besitzt jeder solche Punkt einen eindeutigen Repräsentanten mit $Z=1$. Damit ergibt sich die Einbettung
\[
    \mathbb{A}^2_K
    \longrightarrow
    \mathbb{P}^2_K,
    \qquad
    (x,y)
    \longmapsto
    [x:y:1].
\]
Die Punkte mit $Z=0$ bilden die Gerade im Unendlichen
\[
    L_\infty
    =
    \{[X:Y:Z]\in\mathbb{P}^2_K \mid ZU=0\}.
\]
Die Punkte mit $Z\neq 0$ entsprechen dagegen den affinen Punkten. Um den Punkt im Unendlichen einer elliptischen Kurve zu bestimmen, wird die affine Weierstraßgleichung
\[
    y^2=x^3+Ax+B
\]
homogenisiert. Für einen projektiven Punkt $[X:Y:Z]$ mit $Z\neq 0$ gilt
\[
    x=\frac{X}{Z},
    \qquad
    y=\frac{Y}{Z}.
\]
Durch Einsetzen ergibt sich
\[
    \left(\frac{Y}{Z}\right)^2
    =
    \left(\frac{X}{Z}\right)^3
    +
    A\frac{X}{Z}
    +
    B.
\]
Nach Multiplikation mit $Z^3$ folgt
\[
    Y^2Z
    =
    X^3+AXZ^2+BZ^3.
\]
Damit erhält man das homogene Polynom
\[
    F(X,Y,Z)
    =
    Y^2Z-X^3-AXZ^2-BZ^3.
\]
Da $F$ homogen vom Grad $3$ ist, gilt
\[
    F(\lambda X,\lambda Y,\lambda Z)
    =\lambda^3F(X,Y,Z)
\]
für jedes $\lambda\in K^\times$. Die Bedingung $F(X,Y,Z)=0$ ist daher unabhängig von der Wahl des Repräsentanten. Für einen Punkt auf der Geraden im Unendlichen gilt $Z=0$. Einsetzen in die projektive Weierstraßgleichung liefert
\[
    X^3=0.
\]
Daraus folgt
\[
    X=0.
\]
Da das Tripel $(0,0,0)$ ausgeschlossen ist, muss $Y\neq 0$ gelten. Wegen der Äquivalenz homogener Koordinaten kann $Y$ auf $1$ normiert werden. Somit ergibt sich
\[
    [0:Y:0]
    =
    [0:1:0].
\]
Die projektive Weierstraßkurve besitzt daher genau einen Punkt auf der Geraden im Unendlichen. Dieser Punkt wird mit
\[
    \mathcal{O}
    =
    [0:1:0]
\]
bezeichnet und dient als neutrales Element der Gruppenstruktur elliptischer Kurven. Abschließend kann überprüft werden, dass der Punkt $\mathcal{O}$ nicht singulär ist.

\begin{definition}[Singulärer Punkt]
    Sei $F\in K[X,Y,Z]$ ein von null verschiedenes homogenes Polynom und
    \[
        C = \{[X:Y:Z]\in\mathbb{P}^2_K \mid F(X,Y,Z)=0\}
    \]
    die dadurch definierte projektive Kurve. Ein Punkt \( P=[a:b:c]\in C\) heißt singulär, wenn
    \[
        \frac{\partial F}{\partial X}(a,b,c)
        =
        \frac{\partial F}{\partial Y}(a,b,c)
        =
        \frac{\partial F}{\partial Z}(a,b,c)
        =
        0.
    \]
    Andernfalls heißt der Punkt nicht singulär.
\end{definition}

Für das homogene Polynom $F$ gelten
\[
    F_X=-3X^2-AZ^2,\qquad
    F_Y=2YZ,\qquad
    F_Z=Y^2-2AXZ-3BZ^2.
\]
Am Repräsentanten $(0,1,0)$ des Punktes $\mathcal{O}$ ergibt sich
\[
    F_X(0,1,0)=0,\qquad
    F_Y(0,1,0)=0,\qquad
    F_Z(0,1,0)=1.
\]
Da nicht alle partiellen Ableitungen verschwinden, ist $\mathcal{O}$ nicht singulär.
\begin{figure}[htbp]
    \centering
    \includegraphics[width=0.35\linewidth]{ell_curve_projective_space.png}
    \caption[Projektiver Abschluss einer elliptischen Kurve]{
        Projektiver Abschluss einer elliptischen Kurve.
        Quelle: übernommen aus \cite[S.~202]{forster15}.
    }
    \label{fig:projektiver-abschluss}
\end{figure}

In der projektiven Ebene schneiden sich je zwei verschiedene Geraden in
genau einem Punkt. Insbesondere schneiden sich zwei parallele affine Geraden nach ihrer projektiven Erweiterung in einem Punkt auf der Geraden im Unendlichen. Seien dazu
\[
    g_1: ax+by+c_1=0 \quad \text{und}\quad g_2: ax+by+c_2=0
\]
zwei verschiedene parallele affine Geraden. Ihre projektiven Erweiterungen
lauten
\[
    aX+bY+c_1Z=0 \quad \text{und}\quad aX+bY+c_2Z=0.
\]
Für einen gemeinsamen Schnittpunkt folgt durch Subtraktion
\[
    (c_1-c_2)Z=0.
\]
Da $c_1\neq c_2$ gilt, folgt $Z=0$. Für $Z=0$ verbleibt die Gleichung \(aX+bY=0\). Der gemeinsame Schnittpunkt ist daher \([X:Y:Z]=[b:-a:0]\). Insbesondere haben vertikale Geraden der Form $x=c$ die projektive Gleichung $X-cZ=0$. Sie verlaufen somit alle durch den Punkt \( [0:1:0]=\mathcal{O}\).

\section{Beweis der Assoziativität}
\label{sec:assoziativitaet}
Der folgende Beweis orientiert sich an \cite[S.~477--486]{zwegers24}. Für den allgemeinen Fall werden zunächst explizite Koordinaten für $(P_1+P_2)+P_3$ hergeleitet. Zusammen mit der bereits nachgewiesenen Kommutativität ergibt sich daraus die Assoziativität. Die von dieser Darstellung ausgeschlossenen Sonderfälle werden anschließend gesondert behandelt.

Im gesamten Abschnitt sei \(E:y^2=x^3+Ax+B\) eine elliptische Kurve über einem Körper $K$ mit $\operatorname{char}(K)\neq 2,3$ und $4A^3+27B^2\neq 0$. Zur Vereinfachung der Notation sei außerdem $f(x)=x^3+Ax+B$.

\begin{theorem}
    Seien \(P_1,P_2,P_3\in E\), dann gilt \((P_1+P_2)+P_3=P_1+(P_2+P_3)\).
\end{theorem}

Doch bevor die Assoziativität bewiesen werden kann, müssen zuerst noch eine Reihe von Hilfsaussagen eingeführt werden.

\begin{lemma}
    \label{lem:dreipunktsumme}
    Seien $P_1,P_2,P_3\in E\setminus{\{\mathcal{O}\}}$ mit $P_3\neq \pm P_1, P_2\neq -P_1, P_2\neq -P_3$ und $P_1+P_2\neq -P_3$. Dann besitzt der Punktes $P_4=(P_1+P_2)+P_3$ die Koordinaten
    \[
        x_4=-x_1+x_2-x_3+\frac{1}{c_2^2}-2\frac{c_1}{c_2} \quad \text{und}\quad y_4 = -y_2-(x_4-x_2)(c_1+x_4c_2)
    \]
    wobei
    \[
        c_1 := \frac{x_1m(P_2,P_3)-x_3m(P_1,P_2)}{x_1-x_3} \quad \text{und}\quad c_2:=\frac{m(P_1,P_2)-m(P_2,P_3)}{x_1-x_3}.
    \]
\end{lemma}

\begin{hilfssatz}
    \label{hilfs:fixpunkt-addition}
    Sei $P_1,P_2\in E$ mit $P_1+P_2=P_1$. Dann gilt $P_2 =\mathcal{O}$.
\end{hilfssatz}

\begin{proof}
    Zunächst sei $P_1=\mathcal{O}$. Dann folgt unmittelbar \(P_2=\mathcal{O}+P_2=P_1+P_2=P_1=\mathcal{O}\).

    Sei nun $P_2=-P_1$. Nach Definition der Punktaddition gilt dann \(P_1+P_2=\mathcal{O}\). Aus der Vorraussetzung $P_1+P_2=P_1$ folgt somit $   P_1=\mathcal{O}$ und daher ebenfall $P_2=-P_1=\mathcal{O}$.

    Es bleibt der Fall \(P_1,P_2\in E\setminus\{\mathcal{O}\}\) mit \(P_2\neq-P_1\). Angenommen es gilt dennoch $P_1+P_2=P_1$. Dann gilt für die Summe $P_1+P_2=(x_3,y_3)$ und damit $(x_1,y_1)=(x_3,y_3)$.
    Nach der Additionsformel folgt somit
    \begin{align*}
        y_1 & =y_3=-y_1-m_{12}(x_3-x_1) \\
            & =-y_1-m_{12}\cdot 0       \\
            & = -y_1                    \\
    \end{align*}
    Da \(\operatorname{char}(K)\neq2\) gilt, folgt daraus \(y_1=0\).
    Aus Gleichung~\eqref{eq:faktorisierung-wohldefiniertheit} folgt
    \begin{align*}
        f(x)-\left(y_1+m_{12}(x-x_1)\right)^2
         & = f(x)-m_{12}^2(x-x_1)^2 \\
         & = (x-x_1)(x-x_2)(x-x_3)  \\
         & = (x-x_1)^2(x-x_2).
    \end{align*}
    Wegen \(y_1=0\) folgt durch Umstellen und Ausklammern
    \[
        f(x)
        =
        (x-x_1)^2
        \left(x-x_2+m_{12}^2\right).
    \]
    Folglich ist $x_1$ eine doppelte Nullstelle von $f$. Da zugleich $y_1=0$ gilt, wäre der Punkt $(x_1,0)$ ein singulärer Punkt der Kurve $E$. Dies widerspricht der Voraussetzung, dass $E$ eine elliptische und damit nichtsinguläre Kurve ist. Folglich kann dieser Fall nicht auftreten.
\end{proof}

\textbf{TODO: Hier weiter machen}

\begin{hilfssatz}
    \label{hilfs:steigungsidentitaet}
    Sei $P_1,P_2,P_3\in E\setminus\{\mathcal{O}\}$ mit $P_3 \neq \pm P_1,P_2\neq -P_1,P_2\neq -P_3$ und $P_1+P_2\neq -P_3$. Außerdem ist $c_2$ definiert wie in Lemma~\ref{lem:dreipunktsumme}. Dann gilt
    \[
        c_2(m_{12}+m_{12,3})=1.
    \]
\end{hilfssatz}

\begin{proof}
    Es gilt wie beim Beweis der Wohldefiniertheit der Addition~\ref{lem:wohldefiniertheit}
    ist nach Gleichung~\eqref{eq:faktorisierung-wohldefiniertheit} bekannt, dass
    \[
        f(x)-\left(y_1+m_{12}(x-x_1)\right)^2
        =(x-x_1)(x-x_2)(x-\hat{x}).
    \]
    Wenn $P_2\neq P_3$, dann gilt
    \[
        c_2=\frac{m(P_1,P_2)-m(P_2,P_3)}{x_1-x_3}=\frac{m_{12}-\frac{y_3-y_2}{x_3-x_2}}{x_1-x_3}=\frac{\frac{m_{12}(x_3-x_2)}{x_3-x_2}-\frac{y_3-y_2}{x_3-x_2}}{x_1-x_3}
        = \frac{-y_3+y_2+m_{12}(x_3-x_2)}{(x_1-x_3)(x_3-x_2)}.
    \]
    Wir nutzen dafür, dass $y_2-y_1=m_{12}(x_2-x_1)$ gilt für die Gerade durch die Punkte $P_1$ und $P_2$. Dies formen wir um zu $y_2=y_1+m_{12}(x_2-x_1)$ und setzen es in die Gleichung ein
    \[
        \frac{-y_3+(y_1+m_{12}(x_2-x_1))+m_{12}(x_3-x_2)}{(x_1-x_3)(x_3-x_2)}
        = \frac{-y_3+y_1+m_{12}(x_3-x_1)}{(x_1-x_3)(x_3-x_2)}
    \]
    Wenn nun $P_1+P_2\neq P_3$ gilt. Wir definieren $(x_1,y_1) + (x_2,y_2) = (\hat{x},\hat{y})$ und rechnen
    \[
        m_{12}+m_{12,3}=m_{12}+\frac{y_3-\hat{x}}{x_3-\hat{x}}=m_{12}+\frac{y_3+y_1+m_{12}(\hat{x}-x_1)}{x_3-\hat{x}}
    \]
    \[
        =\frac{m_{12}(x_3-\hat{x})}{x_3-\hat{x}}+\frac{y_3+y_1+m_{12}(\hat{x}-x_1)}{x_3-\hat{x}}
    \]
    \[
        =\frac{y_3+y_1+m_{12}(x_3-x_1)}{x_3-\hat{x}}
    \]
    Wir betrachten nun 4 Fälle.
    \begin{enumerate}
        \item Für $P_2\neq P_3$ und $P_1+P_2\neq P_3$
              \begin{align*}
                  c_2(m_{12}+m_{12,3}) = & \frac{-y_3+y_1+m_{12}(x_3-x_1)}{(x_1-x_3)(x_3-x_2)} \frac{y_3+y_1+m_{12}(x_3-x_1)}{x_3-\hat{x}} \\
                  =                      & \frac{(-y_3+y_1+m_{12}(x_3-x_1))(y_3+y_1+m_{12}(x_3-x_1))}{(x_1-x_3)(x_3-x_2)(x_3-\hat{x})}     \\
                  =                      & \frac{-y_3^2+m_{12}^2(x_3-x_1)^2+2 m_{12}(x_3-x_1)y_1+ y_1^2}{(x_1-x_3)(x_3-x_2)(x_3-\hat{x})}  \\
                  =                      & \frac{-y_3^2+(y_1+ m_{12}(x_3-x_1))^2}{(x_1-x_3)(x_3-x_2)(x_3-\hat{x})}
              \end{align*}
              Es gilt $y_3^2=x^3+Ax+B$ mit $x=x_3$, dann $y_3^2=x_3^3+Ax_3+B$. Mit ~\eqref{eq:faktorisierung-wohldefiniertheit} gilt dann
              \begin{align*}
                  c_2(m_{12}+m_{12,3})= & \frac{-(f(x_3)-(y_1+m_{12}(x_3-x_1))^2)}{(x_1-x_3)(x_3-x_2)(x_3-\hat{x})}   \\
                  =                     & \frac{-(x_3-x_1)(x_3-x_2)(x_3-\hat{x})}{-(x_3-x_)(x_3-x_2)(x_3-\hat{x})}= 1
              \end{align*}

        \item Sei nun $P_2\neq P_3$ und $P_1+P_2=P_3$ es gilt erneut $(\hat{x},\hat{y})=(x_3,y_3)$ mit $y_3=-y_1-m_{12}(x_3-x_1)$
              \[
                  c_2 = \frac{-y_3+y_2+m_{12}(x_3-x_2)}{(x_1-x_3)(x_3-x_2)} = \frac{-2y_3}{(x_1-x_3)(x_3-x_2)}
              \]
              Wir leiten nun die Annahme von ~\eqref{eq:faktorisierung-wohldefiniertheit} ab und setzen dann $x=x_3=\hat{x}$
              \[
                  f'(x)-2m_{12}(y_1+m_{12}(x-x_1))=f'(x_3)-2m_{12}(y_1+m_{12}(x_3-x_1))
              \]
              Es gilt dann $\hat{y}=-y_1-m_{12}(\hat{x}-x_1)$ und $\hat{y}=y_3$ und damit $y_1+m_{12}(x_3-x_1)=-y_3$
              \[
                  f'(x_3)+2m_{12}y_3
              \]
              Recht Seite abgeleitet ergibt
              \[
                  (x-x_2)(x-\hat{x})+(x-x_1)(x-\hat{x})+(x-x_1)(x-x_2)
              \]
              wieder $x=\hat{x}=x_3$ einsetzen
              \[
                  (x_3-x_2)(\hat{x}-\hat{x})+(x_3-x_1)(\hat{x}-\hat{x})+(x_3-x_1)(\hat{x}-x_2)
                  = (x_3-x_1)(\hat{x}-x_2)
              \]
              Also Gesamt ergibt sich
              \[
                  f'(x_3)+2m_{12}y_3 = (x_3-x_1)(\hat{x}-x_2)
              \]
              umformen für $y_3\neq 0$ gilt aus der Bedingung von Wohldefiniertheit
              \[
                  \frac{f'(x_3)}{2y_3}+m_{12} = \frac{(x_3-x_1)(\hat{x}-x_2)}{2y_3}
              \]
              Es ist $m_{12,3}=\frac{f'(x_3)}{2y_3}$ und daraus folgt
              \[
                  m_{12}+m_{12,3}= \frac{(x_3-x_1)(\hat{x}-x_2)}{2y_3}
              \]
              Damit gilt dann offensichtlich mit der Vertauschung des Vorzeichens
              \[
                  c_2(m_{12}+m_{12,3})= \frac{2y_3}{(x_3-x_1)(x_3-x_2)}\frac{(x_3-x_1)(x_3-x_2)}{2y_3}= 1
              \]
        \item Für den Fall $P_2=P_3$ und $P_1+P_2\neq P_3$ haben wir $(x_2,y_2)=(x_3,y_3)$ und dadurch $y_3=y_1+m_{12}(x_3-x_1)$ und
              \[
                  m_{12}+m_{12,3}=\frac{y_3+y_1+m_{12}(x_3-x_1)}{x_3-\hat{x}}= \frac{2y_3}{x_3-\hat{x}}
              \]
              Wir nehmen wieder die Ableitung von der Gleichung von oben und setzen jetzt $x=x_2=x_3$
              \[
                  f'(x_3)+2m_{12}y_3 = (x_3-x_1)(x_3-\hat{x})
              \]
              Dass können wir nun umformen wieder zu
              \begin{align*}
                  m(P_2,P_3)-m(P_1,P_2)                       & =\frac{(x_3-x_1)(x_3-\hat{x})}{2y_3} \\
                  \Leftrightarrow
                  \frac{m(P_2,P_3)-m(P_1,P_2)}{(x_3-\hat{x})} & = \frac{x_3-x_1}{2y_3}               \\ \Leftrightarrow
                  c_2                                         & = \frac{x_3-x_1}{2y_3}
              \end{align*}
              und damit gilt dann für $y_2=y_3\neq 0$ \(c_2(m_{12}+m_{12,3})=1\)
        \item Der Fall $P_2=P_3$ und $P_1+P_2=P_3$ kann nach Lemma 2.1 nicht auftretten
    \end{enumerate}
\end{proof}

Wir haben nun die Hilfsmittel, um Lemma~\ref{lem:dreipunktsumme} zu beweisen.
\begin{proof}
    Seien $P_1,P_2,P_3\in E$ wie in
    Lemma~\ref{lem:dreipunktsumme} vorausgesetzt und sei
    \[
        P_4=(P_1+P_2)+P_3=(x_4,y_4).
    \]
    Weiter sei $P_1+P_2=(\hat{x},\hat{y})$. Nach der Additionsformel gilt
    $\hat{x}=-x_1-x_2+m_{12}^2$ und damit
    \[
        x_4=-\hat{x}-x_3+m_{12,3}^2
        =x_1+x_2-x_3+m_{12,3}^2-m_{12}^2.
    \]
    Nach Hilfssatz~\ref{hilfs:steigungsidentitaet} gilt
    $m_{12,3}=\frac{1}{c_2}-m_{12}$, also
    \[
        \begin{aligned}
            x_4
             & =x_1+x_2-x_3+\left(\frac{1}{c_2}-m_{12}\right)^2-m_{12}^2 \\
             & =x_1+x_2-x_3+\frac{1}{c_2^2}-2\frac{m_{12}}{c_2}.
        \end{aligned}
    \]
    Vergleicht man diesen Ausdruck mit der in
    Lemma~\ref{lem:dreipunktsumme} behaupteten Formel für $x_4$, so ist es ausreichend, $m_{12}$ in der Form $m_{12}=c_1+x_1c_2$ darzustellen. Dies folgt direkt aus den Definitionen von $c_1$ und $c_2$, denn
    \[
        \begin{aligned}
            c_1+x_1c_2
             & =\frac{x_1m_{23}-x_3m_{12}}{x_1-x_3}
            +x_1\frac{m_{12}-m_{23}}{x_1-x_3}       \\
             & =\frac{x_1m_{12}-x_3m_{12}}{x_1-x_3}
            =m_{12}.
        \end{aligned}
    \]
    Damit folgt
    \[
        \begin{aligned}
            x_4
             & =x_1+x_2-x_3+\frac{1}{c_2^2}
            -2\frac{c_1+x_1c_2}{c_2}         \\
             & =-x_1+x_2-x_3+\frac{1}{c_2^2}
            -2\frac{c_1}{c_2}.
        \end{aligned}
    \]
    Um den Ausdruck für $y_4$ zu zeigen, wird zunächst die Additionsformel
    \[
        y_4=-y_3-m_{12,3}(x_4-x_3)
    \]
    verwendet. Für die Punkte $P_2$ und $P_3$ gilt
    \[
        y_3-y_2=m_{23}(x_3-x_2),
    \]
    und damit
    \[
        y_3=y_2+m_{23}(x_3-x_2).
    \]
    Weiter gilt
    \begin{align*}
        m_{23}
         & =\frac{m_{23}(x_1-x_3)}{x_1-x_3}                        \\
         & =\frac{x_1m_{23}-x_3m_{12}+x_3(m_{12}-m_{23})}{x_1-x_3} \\
         & =\frac{x_1m_{23}-x_3m_{12}}{x_1-x_3}
        +x_3\frac{m_{12}-m_{23}}{x_1-x_3}                          \\
         & =c_1+x_3c_2.
    \end{align*}
    Damit folgt
    \[
        y_3=y_2+(x_3-x_2)(c_1+x_3c_2).
    \]

    Die zu zeigende Gleichung
    \[
        y_4=-y_2-(x_4-x_2)(c_1+x_4c_2)
    \]
    ist äquivalent zu
    \[
        y_4+y_2+(x_4-x_2)(c_1+x_4c_2)=0.
    \]
    Nun werden die oben erhaltenen Gleichungen eingesetzt
    \begin{align*}
        y_4=-y_3-m_{12,3}(x_4-x_3) = & -(y_2+(x_3-x_2)(c_1+x_3c_2))-m_{12,3}(x_4-x_3) \\
        =                            & -y_2-(x_3-x_2)(c_1+x_3c_2)-m_{12,3}(x_4-x_3)   \\
    \end{align*}
    Es bleibt somit zu zeigen, dass
    \[
        \begin{aligned}
             & (x_4-x_2)(c_1+x_4c_2)
            -(x_3-x_2)(c_1+x_3c_2)                         \\
             & =(x_4-x_3)c_1
            +\bigl((x_4-x_2)x_4-(x_3-x_2)x_3\bigr)c_2      \\
             & =(x_4-x_3)\bigl(c_1+(x_4+x_3-x_2)c_2\bigr).
        \end{aligned}
    \]
    Da aus der Definition von $x_4$
    \[
        x_4+x_3-x_2
        =x_1+m_{12,3}^2-m_{12}^2
    \]
    und $c_1+x_1c_2=m_{12}$ gilt, folgt
    \[
        \begin{aligned}
             & c_1+(x_4+x_3-x_2)c_2                                               \\
             & =m_{12}+(m_{12,3}^2-m_{12}^2)c_2                                   \\
             & =m_{12}
            +c_2(m_{12,3}-m_{12})(m_{12,3}+m_{12})                                \\
             & \overset{\ref{hilfs:steigungsidentitaet}}{=}m_{12}+m_{12,3}-m_{12}
            =m_{12,3},
        \end{aligned}
    \]
    wobei im vorletzten Schritt
    Hilfssatz~\ref{hilfs:steigungsidentitaet} verwendet wurde. Damit gilt
    \[
        \begin{aligned}
             & (x_4-x_2)(c_1+x_4c_2)
            -(x_3-x_2)(c_1+x_3c_2)   \\
             & =m_{12,3}(x_4-x_3).
        \end{aligned}
    \]
    Daraus lässt sich sich die Additionsformel umformen
    \begin{align*}
        y_4
         & = -y_2-(x_3-x_2)(c_1+x_3c_2)
        -m_{12,3}(x_4-x_3)               \\
         & = -y_2-(x_3-x_2)(c_1+x_3c_2)  \\
         & \quad
        -\Bigl((x_4-x_2)(c_1+x_4c_2)
        -(x_3-x_2)(c_1+x_3c_2)\Bigr)     \\
         & = -y_2-(x_4-x_2)(c_1+x_4c_2).
    \end{align*}
    womit auch die behauptete Formel für $y_4$ gezeigt ist.
\end{proof}

\begin{lemma}
    \label{lem:ass-hilfseigenschaften}
    Für alle $P,P_1,P_2\in E$ gilt,
    \begin{enumerate}[label=(\alph*)]
        \item \(-(-P)=P\)
        \item $-(P_1+P_2)=(-P_1)+(-P_2)$
        \item $(P_1+P_2)+(-P_1)=P_1+(P_2+(-P_1))=P_2$
        \item $P_1+P_2=\mathcal{O} \Rightarrow P_2=-P_1$
    \end{enumerate}
\end{lemma}

\begin{proof}
    \begin{enumerate}[label=(\alph*)]
        \item Sei $P=(x,y)\in E\setminus\{\mathcal{O}\}$, dann ist $-P=(x,-y)\in E\setminus\{\mathcal{O}\}$, also gilt $-(-P)=-(x,-y)=(x,y)=P$. Es gilt nach Definition der Identität $-\mathcal{O}=\mathcal{O}$, damit dann auch \(-(-\mathcal{O})=\mathcal{O}\)
        \item Wenn $P_1 =\mathcal{O}\lor P_2 =\mathcal{O}$ folgt direkt,
              \begin{align*}
                  P_1 =\mathcal{O}: \quad -(P_1+P_2) & =-(\mathcal{O}+P_2)=-P_2=\mathcal{O}+(-P_2) \\
                                                     & = (-\mathcal{O})+(-P_2)=(-P_1)+(-P_2)
              \end{align*}
              \begin{align*}
                  P_2 =\mathcal{O}: \quad -(P_1+P_2) & =-(P_1+\mathcal{O})=-P_1= (-P_1)+\mathcal{O} \\
                                                     & = (-P_1)+(-\mathcal{O})=(-P_1)+(-P_2)
              \end{align*}
              Nun sollen $P_1,P_2\in E \setminus\{\mathcal{O}\}$
              Wir unterscheiden 2 Fälle:
              Erster Fall $P_2=-P_1$, dann gilt $-P_2=-(-P_1)=P_1$ und damit auch
              \[
                  -(P_1+P_2)=-(P_1+(-P_1))=-\mathcal{O}=-\mathcal{O}= (-P_1)+P_1=(-P_1)+(-P_2).
              \]
              Für den anderen Fall gilt \(P_2\neq -P_1\), dann gilt per Definition der Wohldefiniertheit~\ref{lem:wohldefiniertheit}
              \[
                  P_1+P_2=(-x_1-x_2+m_{12}^2,-y_1-m_{12}(x_3-x_1)).
              \]
              Weiter gilt dann auch \(-P_2\neq-(-P_1)\) und \(-P_1=(x_1,-y_1),-P_2=(x_2,-y_2)\), dadurch gilt dann \(m(-P_1,-P_2)=-m(P_1,P_2)\)
              \begin{align*}
                  (-P_1)+(-P_2)= & (-x_1-x_2+m(P_1,P_2)^2,y_1-m(P_1,P_2)(x_3-x_1))
                  =              & -(P_1+P_2).
              \end{align*}
        \item Für $P_1=\mathcal{O}$ oder $P_2=\mathcal{O}$ folgt die Behauptung unmittelbar. Sei also $P_1,P_2\in E\setminus\{\mathcal{O}\}$. Für $P_2=-P_1$ gilt $(P_1+P_2)+(-P_1)=\mathcal{O}+P_2=P_2$. Der Fall $P_1+P_2=P_1$ kann nach Hilfssatz~\ref{hilfs:fixpunkt-addition} nicht auftreten. Sei daher $P_2\neq -P_1$ und $P_1+P_2\neq P_1$ und setze $P_3=-P_1$. Dann gilt $x_3=x_1$ und $y_3=-y_1$.

              Falls $P_1+P_2\neq P_3$, folgt wie in Hilfssatz~\ref{hilfs:steigungsidentitaet}
              \[
                  m_{12}+m_{12,3}
                  =
                  \frac{y_3+y_1+m_{12}(x_3-x_1)}{x_3-\hat{x}}
                  =0.
              \]
              Falls $P_1+P_2=P_3$, gilt entsprechend
              \[
                  m_{12}+m_{12,3}
                  =
                  \frac{(x_3-x_1)(x_3-x_2)}{2y_3}
                  =0,
              \]
              da $x_3=x_1$. Somit gilt stets $m_{12,3}=-m_{12}$.

              Mit $P_1+P_2=(\hat{x},\hat{y})$ und $\hat{x}=m_{12}^2-x_1-x_2$ folgt
              \[
                  x_4=m_{12,3}^2-\hat{x}-x_3
                  =m_{12}^2-(m_{12}^2-x_1-x_2)-x_1=x_2,
              \]
              sowie
              \[
                  y_4=-y_3-m_{12,3}(x_4-x_3)
                  =y_1+m_{12}(x_2-x_1)=y_2.
              \]
              Damit ist $P_4=P_2$ und somit $(P_1+P_2)+(-P_1)=P_2$. Ersetzen wir $P_1$ durch $-P_1$, folgt mit Teil (a) und der Kommutativität
              \[
                  P_2=((-P_1)+P_2)+(-(-P_1))
                  =P_1+(P_2+(-P_1)).
              \]
              Damit gilt
              \[
                  (P_1+P_2)+(-P_1)=P_1+(P_2+(-P_1))=P_2.
              \]
        \item Wenn $P_1+P_2=\mathcal{O}$, dann nutzen wir (c)
              \[
                  P_2=(P_1+P_2)+(-P_1)=\mathcal{O}+(-P_1)=-P_1.
              \]
    \end{enumerate}
\end{proof}

Nun können wir den Beweis der Assoziativität durchführen.
\begin{proof}[Beweis der Assoziativität]
    Wenn $P_1,P_2$ oder $P_3$ gleich $\mathcal{O}$ gilt, dann folgt direkt
    \[
        P_1=\mathcal{O}: \quad (P_1+P_2)+P_3=(\mathcal{O}+P_2)+P_3=P_2+P_3 = \mathcal{O}+(P_2+P_3)=P_1+(P_2+P_3)
    \]
    \[
        P_2=\mathcal{O}: \quad (P_1+P_2)+P_3=(P_1+\mathcal{O})+P_3 =P_1+P_3 = P_1+(\mathcal{O}+P_3)=P_1+(P_2+P_3)
    \]
    \[
        P_3=\mathcal{O}: \quad (P_1+P_2)+P_3=(P_1+P_2)+\mathcal{O} =P_1+P_2 = P_1+(P_2+\mathcal{O})=P_1+(P_2+P_3)
    \]
    Wenn nun $P_3=P_1$, dann gilt aufgrund der Kommutativität der Addition
    \[
        (P_1+P_2)+P_3=(P_1+P_2)+P_1=P_1+P_1+(P_1+P_2)=P_1+(P_2+P_3).
    \]
    Wenn nun $P_3=-P_1$, dann gilt mit der Theorem \ref{lem:ass-hilfseigenschaften} (c)
    \[
        (P_1+P_2)+P_3 = (P_1+P_2)+(-P_1)=P_1+(P_2+(-P_1))=P_1+(P_2+P_3)
    \]
    Für den Fall, dass $P_2=-P_1$ oder $P_2=-P_3$ könenn wir auch \ref{lem:ass-hilfseigenschaften} (c) nutzen
    \begin{align*}
        P_2=-P_3:\quad
        P_2=-P_1:\quad & (P_1+P_2)+P_3
        = (P_1+(-P_1))+P_3=\mathcal{O}+P_3=P_3                                       \\
                       & = P_1+(P_3+(-P_1))=P_1+((-P_1)+P_3)=P_1+(P_2+P_3)           \\
        P_2=-P_3:\quad & (P_1+P_2)+P_3
        = (P_1+(-P_3))+P_3=P_3+(P_1+(-P_3))                                          \\
                       & = P_1 = P_1+\mathcal{O}=P_3=P_1+((-P_3)+P_3)=P_1+(P_2+P_3).
    \end{align*}
    Als nächsten Fall wird $P_1+P_2=-P_3$ betrachtet, durch Anwenden von den Hilfseigenschaften~\ref{lem:ass-hilfseigenschaften} kann man herausfinden, dass
    \[
        P_2+P_3\overset{(a)}{=}P_2+(-(-P_3))=P_2+(-(P_1+P_2))\overset{(b)}{=}P_2+((-P_1)+(-P_2))=-P_1
    \]
    und kann gefolgert werden
    \[
        (P_1+P_2)+P_3=(-P_3)+P_3=\mathcal{O}=P_1+(-P_1)=P_1+(P_2+P_3).
    \]
    Für den letzten Fall wird angenommen, dass $P_1,P_2,P_3\in E\setminus{\mathcal{O}}$ mit $P_3\neq\pm P_1$, $P_2\neq -P_1$, $P_2\neq-P_3$ und $P_1+P_2\neq-P_3$. Das stimmt genau mit den Bedingungen für Lemma~\ref{lem:dreipunktsumme} überein. Zuvor wurde gezeigt, dass
    \[
        P_1+P_2=-P_3
        \quad\Longleftrightarrow\quad
        P_2+P_3=-P_1.
    \]
    Somit folgt aus $P_1+P_2\neq-P_3$ auch
    $P_2+P_3\neq-P_1$, sodass die Voraussetzungen von
    Lemma~\ref{lem:dreipunktsumme} auch nach Vertauschen von
    $P_1$ und $P_3$ erfüllt sind.  Da die dort erhaltenen Koordinaten
    von $P_4$ unter dieser Vertauschung unverändert bleiben, gilt
    \[
        (P_1+P_2)+P_3=(P_3+P_2)+P_1.
    \]
    Mit der Kommutativität der Punktaddition folgt schließlich
    \[
        (P_1+P_2)+P_3
        =(P_3+P_2)+P_1
        =P_1+(P_2+P_3).
    \]
\end{proof}

\begin{bemerkung}[Geometrische Interpretation]
    Sei $K=\mathbb{R}$ und $P_1,P_2\in E\setminus\{\mathcal{O}\}$ mit
    $P_2\neq -P_1$. Die Gerade durch $P_1$ und $P_2$ wird durch
    \[
        l:y=y_1+m_{12}(x-x_1)
    \]
    beschrieben. Für $P_1\neq P_2$ handelt es sich um eine Sekante der
    elliptischen Kurve, für $P_1=P_2$ um die Tangente im Punkt $P_1$.
    In beiden Fällen besitzt die Gerade einen dritten Schnittpunkt mit
    $E$. Dessen Spiegelung an der $x$-Achse ergibt den Summenpunkt
    $P_1+P_2$. Für die Betrachtung der Summe dreier Punkte wird diese geometrische
    Idee erweitert. Anstelle einer Geraden wird eine Parabel
    \[
        p(x)=y_2+(x-x_2)(c_1+xc_2).
    \]
    betrachtet. Das entspricht den Größen $c_1$ und $c_2$ aus Lemma~\ref{lem:dreipunktsumme} .
    Aus den bereits gezeigt Beziehungen
    \[
        c_1+x_1c_2=m_{12}\qquad\text{und}\qquad
        c_1+x_3c_2=m_{23}
    \]
    folgt
    \[
        p(x_1)=y_1, \quad p(x_2)=y_2,\quad
        p(x_3)=y_3.
    \]
    Die Parabel verläuft somit durch die drei Punkte $P_1,P_2$ und $P_3$.
    Sei $K=\mathbb{R}$ und $P_1,P_2\in E\setminus\{\mathcal{O}\}$ mit $P_2\neq -P_1$.
    Es wird nun eine Gerade aufgestellt durch die Punkte $P_1$ und $P_2$
    \[
        l:y=y_1+m_{12}(x-x_1).
    \]
    Nun gibt es verschiedene Fälle. Für $P_2=P_1$ ist die Gerade eine Tangente zur elliptischen Kurve $E$, die sie im Punkt $P_1$ schneidet. Die Tangente hat dann einen dritten Schnittpunkt $(x_3,-y_3)$
\end{bemerkung}

\section{Elliptische Kurven in Charakteristik 2} (nach Rosen, S. 47ff.)
\label{charakteristik-2}
\textbf{Muss man hier auch characteristik 3 ergänzen?}

Computer arbeiten tendenziell besser über Körpern der Charakteristik 2 im Vergleich zu Körpern mit großen prim Charakteristiken. Daher werden häufig Elliptische Kurven über einem Feld $\mathbb{F}_q$ mit $q=2^k$ Elementen verwendet.
(vgl. \cite{silverman06}, S.47)

Aus dieser Notwendigkeit betrachten wir nun genauer Elliptische Kurven in Charakteristik 2, da hier einige unserer vorher definierten Formeln nicht funktionieren.\\
\begin{satz}
    In Charakteristik \(2\) ist die kurze Weierstrass Gleichung \(y^2=x^3+Ax+B\) stets singulär. Daher verwendet man in Charakteristik 2 die allgemeine Weierstraß-Gleichung
\end{satz}

\[
    y^2+a_1xy+a_3y=x^3+a_2x^2+a_4x+a_6.
\]
\begin{proof}
    Betrachte
    \[
        f(x,y) = y^2-x^3-Ax-B.
    \]
    Ein Punkt auf der Kurve ist genau dann singulär, wenn beide partiellen Ableitungen verschwinden. Es gilt
    \[
        f_y=2y\equiv0,
    \] da der Körper Charakteristik 2 hat. Außerdem ist
    \[
        f_x=-3x^2-A
    \]
    Sei nun $x_0$ eine Nullstelle von $f_x$. Wir wählen ein \(y_0\) mit
    \[
        y_0^2=x_0^3+Ax_0+B
    \]
    Dann liegt $(x_0,y_0)$ auf der Kurve und es ist
    \[
        f_x(x_0,y_0)=f_y(x_0,y_0) = 0
    \]
    Damit folgt die Kurve ist singulär.
\end{proof}

Nun wollen auch für Charakterstik 2 die Gruppengesetze aufstellen.
Hierfür beginnen wir mit der allgmeinen Weierstraß-Gleichung für eine elliptische Kurv $E$
\[
    y^2+a_1xy+a_3y=x^3+a_2x^2+a_4x+a_6.
\]
Man schreibt diese in die Form
\[
    f(x,y)=y^2+a_1xy+a_3y-x^3-a_2x^2-a_4x-a_6.
\]
Dann sind die partiellen Ableitungen gegeben durch
\[
    f_x=a_1y-3x^2-2a_2x-a_4,\qquad f_y=2y+a_1x+a_3.
\]
In Charakteristik \(2\) vereinfacht sich dies zu
\[
    f_x=a_1y+x^2+a_4,\qquad f_y=a_1x+a_3.
\]
Ein Punkt \((x,y)\) ist wie bereits früher eingeführt, genau dann singulär, wenn
\[
    f(x,y)=0,\qquad f_x(x,y)=0,\qquad f_y(x,y)=0.
\]
Fall 1: $a_1\neq 0$
Man kann aus \(a_1x+a_3=0\) folgern
\[
    x=\frac{a_3}{a_1}.
\]
Eingesetzt in \(a_1y+x^2+a_4=0\) liefert es,
\[
    a_1y+\left(\frac{a_3}{a_1}\right)^2+a_4=0
\]
\[
    y= \frac{a_3^2}{a_1^3}+ \frac{a_4}{a_1}.
\]
wenn der Punkt zusätzlich \(f(x,y)=0\) erfüllt handelt es sich um einen singulären Punkt.

Fall 2: $a_1=0$
Es gilt nun,
\[
    f_x=x^2+a_4,\qquad f_y=a_3.
\]
also muss \(a_3= 0\) und $f(x,y)=0$ gelten, sodass $(x,y)$ ein singulärer Punkt ist.\\

Nun können wir die Addition von Punkten definieren. Es gilt wieder für alle Punkte $P\in E$
\[
    P+\mathcal{O}=P.
\]
\textbf{washington S.48-49 hier weiter machen}

